%------------------------------------------------------------------------------%

\usepackage{integracao_e_series}

\title{Séries Numéricas -- Teste da Razão}

%------------------------------------------------------------------------------%
\begin{document}

\addtitle

%------------------------------------------------------------------------------%
%------------------------------------------------------------------------------%
\section{Teste da Razão}

%------------------------------------------------------------------------------%
\begin{frame}{Motivação -- Analise da Série Geométrica}
	Considere a série geométrica com $a$ e $r$ positivos
	\[
    \sum_{n=1}^{\infty} ar^{n-1}
  \]
	\pause
	Podemos comparar o crescimento de seus termos calculando a razão
	\[
    \rho
    = \frac{\,a_{n+1}\,}{a_n}
    \pause
    = \frac{ar^n}{\,ar^{n-1}\,}
    \pause
    = r
  \]

	\pause
  \vspace{0.5\baselineskip}
  Se $\rho<1$     \tabto{17mm} a série converge

  \pause
  \vspace{\baselineskip}
  Se $\rho\geq 1$ \tabto{17mm} a série diverge
\end{frame}

%------------------------------------------------------------------------------%
\begin{frame}{Teste da Razão}
	Seja $\;\sum a_n\;$ tal que $\;a_n>0\;$ para $\;n\geq N\;$ \pause e
	\[
     \rho = \lim_{n\to\infty}\frac{a_{n+1}}{a_n}
  \]

  \pause
  \vspace{\baselineskip}
	Se $\rho<1$                    \tabto{37mm} a série \emph{converge}

  \pause
  \vspace{\baselineskip}
	Se $\rho>1$ ou $\rho\to\infty$ \tabto{37mm} a série \emph{diverge}

  \pause
  \vspace{\baselineskip}
	Se $\rho=1$                    \tabto{37mm} o teste é \emph{inconclusivo}
\end{frame}

%------------------------------------------------------------------------------%
\addtocounter{exemplo}{1}
\begin{frame}{Exemplo \theexemplo}
	\onslide<+->{Analisar a convergência da série
	\[
    \sum_{n=1}^{\infty} \frac{(2n)!}{(n!)^2}
  \]}

	\onslide<+->{Calculando os termos}
	\begin{align*}
		\onslide<2->{a_n & = \frac{(2n)!}{n!n!} \\[2mm]}
		\onslide<+->{a_{n+1} & = \frac{\big(2(n+1)\big)!}{(n+1)!(n+1)!}}
	\end{align*}
\end{frame}

%------------------------------------------------------------------------------%
\begin{frame}{Exemplo \theexemplo}
	\begin{align*}
		\onslide<+->{\frac{a_{n+1}}{a_n}}
		\onslide<+->{& = \left(\,\frac{(2n+2)!}{\,(n+1)!(n+1)!\,}\,\right) \; \left(\,\frac{n!n!}{\,(2n)!\,}\,\right) \\[2mm]}
		\onslide<+->{& = \frac{n!}{(n+1)!} \; \frac{n!}{(n+1)!} \; \frac{(2n+2)!}{(2n)!}                              \\[2mm]}
		\onslide<+->{& = \frac{n!}{(n+1)n!} \; \frac{n!}{(n+1)n!} \; \frac{(2n+2)(2n+1)(2n)!}{(2n)!}                  \\[2mm]}
		\onslide<+->{& = \frac{1}{(n+1)} \; \frac{1}{(n+1)} \; (2n+2)(2n+1)                                           \\[2mm]}
		\onslide<+->{& = \frac{\,2(n+1)(2n+1)\,}{(n+1)(n+1)}}
		\onslide<+->{\; = \;\frac{\,4n+2\,}{n+1}}
		\onslide<+->{\; \to\; 4}
    \onslide<+->{\; > 1}
		\onslide<+->{\qquad\text{série divergente}}
	\end{align*}
\end{frame}

%------------------------------------------------------------------------------%
%------------------------------------------------------------------------------%
\section{Exemplos do Teste da Razão}

%------------------------------------------------------------------------------%
\addtocounter{exemplo}{1}
\begin{frame}{Exemplo \theexemplo}
  Use o teste da razão para verificar a convergência ou divergência da série
  \[
    \sum_{n=1}^\infty\frac{3^{n+2}}{\ln(n)}
  \]
\end{frame}

%------------------------------------------------------------------------------%
\begin{frame}{Exemplo \theexemplo}
  \[
    a_n = \frac{\,3^{n+2}\,}{\,\ln(n)\,}
  \]
  \pause
  \[
    \frac{a_{n+1}}{a_n}
     = \frac{3^{n+3}}{\,\ln(n+1)\,} \; \frac{\,\ln(n)\,}{\,3^{n+2}\,}
     \pause
     = \frac{3\ln(n)}{\,\ln(n+1)\,}
  \]

  \pause
  \vspace{0.5\baselineskip}
  \[
    \rho
    = \lim_{n\to\infty} \frac{\,a_{n+1}\,}{a_n}
    \pause
    = \lim_{n\to\infty} \frac{3\ln(n)}{\,\ln(n+1)\,}
    \pause
    = \lim_{x\to\infty} \frac{3\ln(x)}{\,\ln(x+1)\,}
    \pause
    = \lim_{x\to\infty} \frac{\,3(x+1)\,}{x}
    \pause
    = 3
    \pause
    > 1
  \]

  A série diverge
\end{frame}

%------------------------------------------------------------------------------%
\addtocounter{exemplo}{1}
\begin{frame}{Exemplo \theexemplo}
  Use o teste da razão para verificar a convergência ou divergência da série
  \[
    \sum_{n=1}^\infty \frac{n!}{\,n^n\,}
  \]
\end{frame}

%------------------------------------------------------------------------------%
\begin{frame}{Exemplo \theexemplo}
\begin{columns}
\column{0.5\linewidth}
  \onslide<+->{\[
    a_n = \frac{n!}{\,n^n\,}
  \]}
  \begin{align*}
    \onslide<+->{\frac{\,a_{n+1}\,}{a_n}
    & = \frac{(n+1)!}{\,(n+1)^{n+1}\,} \; \frac{\,n^n\,}{n!} \\[3mm]}
    \onslide<+->{& = \frac{\,(n+1)\,n!\;n^n\,}{(n+1)^n\,(n+1)\;n!}        \\[3mm]}
    \onslide<+->{& = \frac{n^n}{\,(n+1)^n\,}}
  \end{align*}
\column{0.5\linewidth}
  \begin{align*}
    \onslide<+->{\frac{\,a_{n+1}\,}{a_n}
    & = \left(\frac{n}{\,n+1\,}\right)^{\!n}                     \\[3mm]}
    \onslide<+->{& = \left(\frac{1}{1+\sfrac{1}{n}}\right)^{\!n}              \\[3mm]}
    \onslide<+->{& = \left[\left( 1 + \frac{1}{n} \right)^{\!n\vphantom{\int}}\;\right]^{-1}}
  \end{align*}
\end{columns}
\end{frame}

%------------------------------------------------------------------------------%
\begin{frame}{Exemplo \theexemplo}
  \begin{align*}
    \onslide<+->{\rho
    & = \lim \frac{a_{n+1}}{a_n} \\[3mm]}
    \onslide<+->{& = \lim \left[\left( 1 + \frac{1}{n} \right)^{\!n\vphantom{\int}}\;\right]^{-1} \\[3mm]}
    \onslide<+->{& = \left[ \lim \left( 1 + \frac{1}{n} \right)^{\!n\vphantom{\int}}\;\right]^{-1} \\[3mm]}
    \onslide<+->{& = e^{-1} }
    \onslide<+->{\approx 0,3679}
    \onslide<+->{< 1}
  \end{align*}
  \onslide<6->{A série converge}
\end{frame}
%------------------------------------------------------------------------------%

%------------------------------------------------------------------------------%
\addtocounter{exemplo}{1}
\begin{frame}{Exemplo \theexemplo}
  Use o teste da razão para verificar a convergência ou divergência da série
  \[
    a_1 = 5,
    \quad
    a_{n+1} = \frac{\,\sqrt[n]{n}\,}{2} a_n
  \]
\end{frame}

%------------------------------------------------------------------------------%
\begin{frame}{Exemplo \theexemplo}
  \[
    a_{n+1} = \frac{\,\sqrt[n]{n}\,}{2} a_n
  \]

  \pause
  \[
    \frac{a_{n+1}}{a_n}
    = \frac{\sqrt[n]{n}}{2}
    \pause
    \to \frac{1}{2}
    \pause
    < 1
  \]
  A série converge
\end{frame}

%------------------------------------------------------------------------------%
\section{Lista Mínima}

%------------------------------------------------------------------------------%
\begin{frame}{Lista Mínima}
  Estudar as Seção 6.6 da Apostila

  \vspace{\baselineskip}
  Exercícios: 3a-f, 4a-f

  \vfill
  Atenção: A prova é baseada no livro, não nas apresentações
\end{frame}

%------------------------------------------------------------------------------%
\end{document}
%------------------------------------------------------------------------------%
